\section{Problem Formulation of Core Loss Modeling}
Let $f$, $B$, and $P_{v}$ denote frequency, peak flux density, and volumetric loss, respectively, and let $\mathcal{D}=\{(f_i, B_i, P_{v,i})\}_{i=1}^{N}$ denote $N$ measured operating points. The focused discovery problem under sinusoidal drive can be expressed as a linear form of nonlinear functions with corresponding coefficients:
\begin{equation}
    P_{v} = \boldsymbol{\Theta}(f, B) \cdot \boldsymbol{\xi}, \quad \boldsymbol{\Theta}(f, B)\in\mathbb{R}^{M}, \quad \boldsymbol{\xi}\in\mathbb{R}^{M}
    \label{eq:ssi_form}
\end{equation}
where the library $\boldsymbol{\Theta}(f, B)$ collects $M$ symbolic candidate functions and $\boldsymbol{\xi}=\left[\xi_1, \xi_2, \cdots, \xi_M \right]^{T}$ is the coefficient vector to be identified. An explicit equation is obtained once $\boldsymbol{\xi}$ is sparse (i.e., most of elements are zero), which selects the governing terms while other zero entries are pruned. The equation to be identified comes from the library, so the design of $\boldsymbol{\Theta}$ determines what the framework is able to discover. Accordingly, the library is constructed to span the reported loss mechanisms while remaining readable in closed form:
\begin{equation}
    \boldsymbol{\Theta}(f, B) =
    \Big[\underbrace{\bar{f}^{\alpha}\bar{B}^{\beta}}_{\text{power laws}} ,
    \underbrace{\tfrac{\bar{f}^{\alpha}\bar{B}^{\beta}}{1+(\bar{f}/f_c)^{2}}}_{\text{HF roll-off}} ,
    \underbrace{\bar{f}^{\alpha}\bar{B}^{\beta}\ln\bar{B},\;\bar{f}^{\alpha}\bar{B}^{\beta}\ln\bar{f}}_{\text{drifting exponents}} ,
    \underbrace{\bar{f}\bar{B}^{2}\big(e^{\bar{B}/B_c}\!-\!1\big)}_{\text{saturation}} ,
    \underbrace{\bar{f}\bar{B},\bar{f},\bar{B},1}_{\text{auxiliaries}} \Big].
    \label{eq:library}
\end{equation}

In the power-law block, the exponents following the Steinmetz-type scaling are designed to be fractional, e.g., $\alpha\in\{1.0,1.2,\cdots,2.2\}$ and $\beta\in\{1.6,1.8,\cdots,3.0\}$, including classical hysteresis, eddy-current and excess contributions and giving $73$ terms. The roll-off block accounts for the non-monotone behavior of $\alpha$ at high frequency in a form $1/[1+(\bar{f}/f_c)^{2}]$ with different exponent pairs, giving $18$ terms. The logarithmic block with $6$ terms is the first-order expansion of an exponent that drifts with the excitation level. The saturation block adds $2$ terms, and $4$ auxiliaries complete the library, giving $M=103$ candidate functions in total. Note that $\alpha$, $\beta$, $f_c$ and $B_c$ are fixed once the grid is selected.


% The roll-off block accounts for the non-monotone behavior of the apparent exponent $\alpha$ at high frequency. Band-wise Steinmetz fits on the present dataset give $\alpha=1.27$, $1.58$ and $1.44$ over the lower, middle and upper frequency thirds, and this turnover cannot be produced by any positive combination of pure power laws. A single-pole attenuation of the form $1/[1+(\bar{f}/f_c)^{2}]$ reproduces it directly, and the cutoffs are placed at $f_c\in\{0.5,0.75,1.0,1.5,2.0,3.0\}$ in units of $f_n$ so that the grid brackets the observed turnover from both sides. Combined with three exponent pairs, this block contributes $18$ terms.

% The logarithmic block is the first-order expansion of an exponent that drifts with the excitation level. Writing the true exponent as $\beta_0+\delta$ and expanding for a small perturbation $\delta$,
% \begin{equation}
%     \bar{B}^{\,\beta_0+\delta} = \bar{B}^{\,\beta_0}e^{\delta\ln\bar{B}}
%     \simeq \bar{B}^{\,\beta_0} + \delta\,\bar{B}^{\,\beta_0}\ln\bar{B}
%     \label{eq:drift}
% \end{equation}
% so that a term $\bar{f}^{\alpha}\bar{B}^{\beta}\ln\bar{B}$ is precisely the leading correction to a drifting $\beta$, and $\bar{f}^{\alpha}\bar{B}^{\beta}\ln\bar{f}$ the corresponding correction to a drifting $\alpha$. This is not a curve-fitting device: on the present dataset $\beta$ measures $2.48$ at low frequency against $2.14$ at high frequency, which is exactly the drift that \eqref{eq:drift} represents. Six such terms are included. The saturation block adds two terms for the loss upturn as $B$ approaches the knee of the $B$--$H$ curve, and four low-order auxiliaries complete the library, giving $M=103$ candidate functions in total.







% \begin{figure}[t]
% \begin{minipage}{0.49\textwidth}
% \subfloat[\label{fig:licking_tiger}A licking tiger]{\includegraphics[width=0.45\textwidth]{figures/tiger1.jpg}} \quad
% \subfloat[\label{fig:mother_tiger}A mother tiger]{\includegraphics[width=0.45\textwidth]{figures/tiger2.jpg}}
% \caption{\label{fig:beautiful_tiger}Beautiful tigers}
% \end{minipage}
% %%\hfill
% \begin{minipage}{0.49\textwidth}
% \subfloat[\label{fig:smiling_tiger}A smiling tiger]{\includegraphics[width=0.45\textwidth]{figures/tiger3.jpg}} \quad
% \subfloat[\label{fig:phd_tiger}A Ph.D. tiger]{\includegraphics[width=0.45\textwidth]{figures/tiger4.jpg}}
% \caption{ \label{fig:more_beautiful_tiger}More beautiful tigers. Everyone loves tigers.}
% \end{minipage}
% \end{figure}
